\scanpage{27}
\uline{Démonstration} \inserted{\textcircled{a}} Par les lemmes 1 \inserted{Cor. 2} et 2, l'équivalence des
conditions (i) \deleted{\unsure{(ii)}} $(\mathrm{ii}_m)$ \deleted{\unsure{(iii)}} est immédiate par récurrence
sur $n$, compte tenu que (i) elles sont ttes
vérifiées pour $n=0$ (\uline{descente fid plate}) (ii) la condition
du lemme 1 \deleted{relatif à} $\mathbb M_m(B^{(n)})=\mathbb M_m(B^{(n)})^*+\operatorname{Im}\mathbb M_m(B^{\unsure{(n+1)|(n)}})$
est satisfaite \deleted{pour} $\forall\,n$, en vertu du lemme 2 (ii)
et de la théorie des relations d'équivalence
artiniennes \ldots. La condition $(\mathrm{ii}_m)$ signifie que $\forall n$,
$\operatorname{Spec} A''_n\unsure{-|\to}\operatorname{Spec} B^{(n)}$ est un morphisme de descente
effective \deleted{\unsure{uni}} pour la catégorie des Modules qu. coh.
\uline{libres} de \deleted{type fini} \inserted{rang $m$}. Cela signifie aussi que tte
donnée de descente sur le $A'$-module $A'_n{}^m$,
relative à $A'_n/B^{(n)}$, cette donnée de descente
est effective dans la catégorie des modules \uline{plats} ;
c'est évidemment néc., et c'est suffisant, car
\inserted{je dis que} le module descendu $F_{B^{(n)}}$ sur $B^{(n)}$ (\deleted{\unsure{p.}}plat sur $B^{(n)}$) est
nécessairement libre de rang $m$ \inserted{(car il le devient après chgt de base $B^{(n)}\to A'_n$)}. Comme $B^{(n)}$ est
\margin{lemme à dégager}
artinien, il suffit de voir que le module
réduit mod $\operatorname{rad}(B^{(n)})$ l'est, ce qui nous ramène
au cas où l'anneau de base est \unsure{une} \inserted{somme finie de} corps, puis
un \deleted{somme de} corps, où c'est immédiat \ldots.

Donc il est clair que $(\mathrm{iii})\Rightarrow(\mathrm{ii})_m$. Pour voir que
% archive sheet number: 26
